% =====================================================================
%  [H. Brøchner], Om Søren Kierkegaards Virksomhed som religieus Forfatter
%  Fædrelandet, 16de Aargang, Løverdagen den 1. December 1855, Nr. 281,
%  printed pp. 1179--1180.  Published anonymously, signed "-r".
%  TRANSCRIPTION (Danish) — diplomatic, from the page images.
% ---------------------------------------------------------------------
%  Scan: ~/bibliotek/Brøchner, Hans/1855-12-01-faedrelandet-nr281.pdf
%        sha256 d3a953dff55bb9bf65cf7fe1977b81c50d0be52b0f0880d501783c7aaf43750f
%        5 PDF pages; PDF 1 is Mediestream's licence sheet and is NOT part of
%        the newspaper.  PAGE MAP: printed p. N = PDF page N − 1177, so
%        1179 = PDF 2 and 1180 = PDF 3.  Downloaded from Mediestream
%        (Det Kgl. Bibliotek) 2026-09-22, record
%        doms_aviser_page:uuid:de038ba0-f5b7-4212-8b78-f52bda29a058.
%
%  TYPE: the body is ANTIQUA (roman) in three columns, NOT Fraktur — the
%        book brief said Fraktur and was wrong.  Emphasis is by
%        letterspacing -> \emph{}.  There is NO italic in the article.
%        One phrase is bold: "gjore opmærksom" (p. 1179).
%
%  EXTENT: the article opens at the head of col. 1 on the front page and
%        runs through all three columns; it continues at the head of col. 1
%        of p. 1180 and ends a sixth of the way down col. 3, at "...i det
%        Liv, han har vakt i saa Manges Sjæle.", followed on the same line
%        by the signature "—r." and, two lines below, a centred rule.
%        Only the article is transcribed; the surrounding newspaper matter
%        is described in comments at the head and foot of each page.
% ---------------------------------------------------------------------
%  THIS FILE REPLACES A WITHDRAWN TEXT.  Until 2026-09-22 this book's
%  transcription had been made from Carl Henrik Koch's 2016 critical
%  edition (Scientia Danica H8 vol. 13) — a modern, critical and
%  in-copyright text, with invented page markers.  REPAIR-PLAYBOOK §8
%  forbids patching derived text against a scan, so it was not patched:
%  it was discarded and the article transcribed afresh from the images,
%  by transcribers who were kept away from both it and Koch's edition.
%  The withdrawn file is kept beside this one as
%  transcription.tex.bak.koch-withdrawn-20260922.
% ---------------------------------------------------------------------
%  ACCURACY  [2026-09-22]
%  Method: each printed page rendered at 400 dpi (600--1200 dpi and the
%  native 200 ppi raster for glyph questions) and transcribed from the
%  image by one agent; then read again, word by word, by a SECOND reader
%  who had not transcribed it; then the disputed sites settled by a THIRD
%  reader measuring the native raster.  The scan's text layer is
%  unusable (SCANS.tsv: WEAK — Mediestream's own phrase search cannot find
%  this article), so ocrdiff was not run and nothing was taken from it.
%
%  Second reading: 0 word-level discrepancies on p. 1179; 1 on p. 1180.
%  Corrections applied after the third reading:
%   * p. 1180 "latterlig Overdrivelse." — a full stop, not a comma (round,
%     6 px, +0.9 below the baseline; the comma of "Alvor," on the same line
%     is 9 px with a tail +4.5 below).  The following "det" is printed
%     lowercase and stands as printed.
%   * p. 1180 "tilbage i Enfold" — ordinary roman, not italic.
%   * p. 1179, 18 emphasis sites — the letterspacing had been marked over
%     whole phrases where the page spaces only one or two words (e.g.
%     "totalt" not "totalt betragtet"; "Troen", "Syndsbevidstheden",
%     "Øieblikket", "Guden i Tiden" but not "et nyt Organ;" or the
%     "en ny ..." phrases).  Measured by inter-letter advance within the
%     line: letterspaced runs give 6--9 px at 400 dpi, plain words 2--5.
%   * two apparatus notes were wrong and are corrected at the site: the
%     footnote quotation „Enkelte, hwem jeg ... min Læser“ DOES close, and
%     the "i" of „Øieblikket“ is a turned sort, not a battered letter.
%
%  FURTHER CORRECTIONS 2026-09-23, found while conforming the translation:
%   * p. 1180 „afmægtig valgte han Intet at gjøre" — the file had
%     "almægtig", a TRANSCRIBER'S error that both readings missed. The
%     second letter carries the f-hook at the top of its stem, matching the
%     f of "fordømme" on the next line; the l of "valgte" on the same line
%     has none. It surfaced because the older translation read "impotent".
%     Word-level errors found in all: 2 in 2 pp.
%   * Three spurious paragraph breaks introduced when the page fragments
%     were spliced, not present in the print: at the col. 1->2 and col. 2->3
%     breaks of p. 1179 (both column heads measured FLUSH at the native
%     raster) and at the page turn, which falls mid-word — now joined as
%     "Reli\opage{1180}gieuse", the printed hyphen resolved.
%
%  OPEN:
%   * ø/o.  The native raster is 200 ppi and a slash is often simply lost.
%     Slashless forms are kept as printed at "Soren" (heading), "gjore",
%     "ifores" and "Loverdagen" (issue line); a slash is supplied at "tør",
%     "udførligere" and "Indøvelse".  Neither treatment is decidable from
%     this scan, and the inconsistency is recorded rather than resolved.
%   * "Guden i Tiden": whether the one-letter word "i" belongs to the
%     letterspaced run cannot be measured; the run is marked whole.
%   * The signature "—r." is an over-inked blob; it is read from its arm,
%     its width and the point, not from a legible letter.
%   * Quote balance is off by design: several quotations are opened in the
%     print and never closed (logged at each site).
% =====================================================================
\documentclass[12pt,a4paper]{book}
\usepackage[utf8]{inputenc}
\usepackage[T1]{fontenc}
\usepackage{libertinus}
\usepackage{libertinust1math}
\usepackage{graphicx}    % for \reflectbox (the old Danish »ɔ:« id-est mark)
\DeclareUnicodeCharacter{0254}{\reflectbox{c}}  % ɔ (U+0254) = reversed c
\usepackage[danish]{babel}
\usepackage[protrusion=true,expansion=true]{microtype}
\usepackage{setspace}
\usepackage[margin=1.2in]{geometry}
\usepackage{fancyhdr}
\setlength{\headheight}{28pt}
\addtolength{\topmargin}{-13.5pt}
\usepackage{hyperref}

\hypersetup{
  pdftitle={Om Søren Kierkegaards Virksomhed som religieus Forfatter},
  pdfauthor={H. Brøchner},
  hidelinks
}

\pagestyle{fancy}
\fancyhf{}
\fancyhead[LE,RO]{\thepage}
\fancyhead[RE]{\textit{Om Søren Kierkegaards Virksomhed}}
\fancyhead[LO]{\textit{Fædrelandet, 1. December 1855}}

\setstretch{1.3}

\title{\textbf{Om Søren Kierkegaards Virksomhed som religieus Forfatter}}
\author{[H.\ Brøchner], signeret „-r“}
\date{Fædrelandet, 16de Aarg., Nr.\ 281, Løverdagen den 1.\ December 1855}

\usepackage{marginnote}
\newcommand{\opage}[1]{\marginnote{\footnotesize\textit{[#1]}}}

\begin{document}
\maketitle
\thispagestyle{empty}

\opage{1179}

% ====================================================================
% Fædrelandet, 16de Aarg., Nr. 281, Løverdagen den 1. December 1855,
% printed p. 1179 = PDF page 2 of
% ~/bibliotek/Brøchner, Hans/1855-12-01-faedrelandet-nr281.pdf
% sha256 d3a953dff55bb9bf65cf7fe1977b81c50d0be52b0f0880d501783c7aaf43750f
% (hash verified against BATCH-AGENT.md before transcribing).
% Folio "— 1179 —" read at the foot of the leaf, centred under col. 2.
%
% WHAT PRECEDES THE ARTICLE ON THIS PAGE: only the masthead apparatus —
%   the "Bladets Comptoir" notice (upper left), the title "Fædrelandet. 1855.",
%   the "Bekjendtgjørelser" advertising-rates notice (upper right), the
%   subscription line, and the issue line "16de Aarg. | Løverdagen den
%   1. December. | Nr. 281". No other article stands on p. 1179.
% WHAT FOLLOWS IT ON THIS PAGE: nothing. The article opens at the head of
%   col. 1 and fills all three columns, breaking off in mid-word at the foot
%   of col. 3 ("Det Reli-") to continue on p. 1180. The only other matter
%   below col. 3 is the article's own footnote *), set under a short rule.
%
% TYPE: the body of this page is ANTIQUA throughout, NOT Fraktur. (The brief
%   anticipated Fraktur; it is not what the page uses. The only display face
%   is the fat titling of the heading.) Emphasis is therefore carried by
%   LETTERSPACING alone, rendered \emph{} below; there is no antiqua/Fraktur
%   contrast to record. One phrase is additionally set in a bold face — noted
%   at the site.
% QUOTATION MARKS: „…“ low-high as printed. Quotations inside quotations
%   occur in col. 3 and are kept apart.
% ====================================================================

\begin{center}
{\large Om Soren Kierkegaards Virksomhed\\
som religieus Forfatter.}
\end{center}
% PRINTED AS IS: "Soren", not "Søren". Checked at 400 and 900 dpi: the o of
% the fat titling face carries no slash, though the same heading's other
% sorts are clean. Transcribed as printed.
% The heading is followed by a short centred rule before the text.

Det har ikke kunnet være anderledes, end at den Mands Død, der hos saa
Mange havde vakt alvorlige Tanker, stærke Lidenskaber, Kjærlighedens og
Hadets, har maattet sætte disse Tanker, disse Lidenskaber, i en levende
Bevægelse. I et saadant Øieblik ligger det nær, at man stræber at overskue
det nu Afsluttede, at man i det klare Billede af hvad han har villet og
virket søger et sikkert Tilhold for sine Tanker og for sin Dom. Til at give
et saadant Billede, som Forfatteren af disse Linier ikke har kunnet finde i
det, der efter Kierkegaards Død er skrevet om ham, skal denne Fremstilling
bidrage, der helt igjennem slutter sig til hans egne Ord, og vil forsøge
dels at vise den indre Sammenhæng og Enhed i hele Rækken af hans Skrifter,
dels hvor nøie hans sidste Optræden hænger sammen med hans tidligere
Virksomhed, hvor bestemt den var antydet i denne, og med hvor nødvendig en
Conseqvens den under de bestemte Forhold maatte fremgaae deraf.

Den Forfattervirksomhed, der her skal være Gjenstand for Omtale, begynder i
Februar 1843 med \emph{Enten---Eller}. Forud for dette Værk havde
Kierkegaard vel udgivet to Skrifter: \emph{Af en endnu Levendes Papirer,
udgivet imod hans Villie af S. K.} (1838), og Magisterdisputatsen:
\emph{Om Begrebet Ironi, med stadigt Hensyn til Socrates} (1841). Men om end
det første af disse Skrifter i Polemiken mod den Enkeltes Opgaaen i den
upersonlige Masse, i Opfattelsen af den foregaaende Generations Forhold til
den nærværende, i Antydningen af det religieuse Maal (S. 30), og endelig i
Meget i Formen (især i Fortalen) minder om haus sildigere Arbeider, saa
% PRINTED AS IS: "haus" for "hans" — a turned n. Verified at 600% on the
% 400 dpi render: the sort is unmistakably u, with the arch at the foot.
staaer det dog ved sin tilfældige Foranledning og begrændsede Gjenstand
isoleret, udenfor den sammenhængende Række; --- og om end det andet og uden
Sammenligning betydeligere Skrift: Afhandlingen om Ironien, udvikler et
Begreb, der for K.s Opfattelse af Existensstadierne er af høi Vigtighed, og
om det end (i 2det Afsnit, navnlig S. 287---304) giver den kategoriske
Bestemmelse af hvad der i Enten---Eller's første Del optræder i personlig
Existens, og i anden Del finder sin Dom, saa er det væsenligen et lærd
Skrift, i hvilket det philologiske og historiske Apparat indtager den
overveiende Plads, og det, der giver det en almindeligere Betydning, er
atter optaget paa sit Sted i de senere Skrifter, navnlig i „Afsluttende
uvidenskabelig Efterskrift“.

Den egenlige, sammenhængende Forfattervirksomhed begynder altsaa med
\emph{Enten---Eller}; den naaer sit afgjørende Punct i „\emph{Afsluttende
uvidenskabelig Efterskrift}“ (Febr. 1846), sin foreløbige Afslutning i
„\emph{To Taler ved Altergangen om Fredagen}“ (Eftersommeren 1851; jfr.
Fortalen til disse) sin endelige først ved Forfatterens Død.
„\emph{Afsluttende Efterskrift}“, der, hvis noget enkelt af K.s Skrifter
skal fremhæves, maa betegnes som hans Hovedværk, danner Midtpunktet i
Rækken af hans Skrifter og deler disse i to Grupper, af hvilke den første
stræber hen imod det Christelige, den anden fremstiller det saaledes naaede
Christelige.

Det Christelige er altsaa Maalet, Opgaven, og er det fra først af.
„Forfatterskabet, \emph{totalt} betragtet, er religieust fra Først til
Sidst“. --- „Problemet, hele Forfatterskabets Problem er: \emph{det at blive
Christen}“. --- „Til dette Forfatterskab svarer som Ophav En, der qua
Forfatter kun har villet Eet“. --- Nærmere er da dette, han har villet,
bestemt saaledes: „\emph{Uden Myudighed} at \textbf{gjore opmærksom} paa det
Religieuse, det Christelige“.
% PRINTED AS IS: "Myudighed" for "Myndighed" — a second turned n, in a
% letterspaced run, so the sort stands clear of its neighbours.
% PRINTED AS IS: "gjore" without the ø-slash. This phrase alone is set in a
% BOLD face inside the letterspaced quotation; checked again at 900 dpi, the
% bold o carries no slash. Rendered \textbf{} to keep the page's two devices
% (letterspacing, bold) apart.
(Om min Forfattervirksomhed, S. 7, 9, 14). Men de, der skulde gjøres
opmærksomme, vare de, der i „Christenhed“ leve i „det Sandsebedrag: at
kalde sig Christen, maaskee indbilde sig at være det, uden at være det“.
Opgaven var da ikke nærmest, \emph{at give}, men først ligesom \emph{at fratage},
at skaffe det bort, der begunstigede Sandsebedraget, og saaledes ved det
Maieutiske, ved denne „Jordemoder-Konst“, som Sokrates kaldte den, at
forhjælpe den Enkelte til selv „at føde Sandheden“. Men Maaden, paa hvilken
der skulde virkes mod Sandsebedraget, var betinget ved dettes Grund. Denne
laa dels i Tidens hele Retning: at den, reflecterende og lidenskabsløs,
havde ladet Forstanden nivellere, og ladet „den Enkelte“ forsvinde i
„Abstractionens Phantom, Publikum“, at den, „væsenlig forstandig, maaskee i
Gjennemsnit vidende som ingen Generation var det før den“, havde „glemt
at \emph{existere} og hvad \emph{Inderlighed} er“, at den,
% ---- COLUMN BREAK: col. 1 ends "…og hvad Inderlighed er“, at den,"
%      col. 2 opens "upersonligt og interesseløst fortabende sig…"
%      (no word is broken across the break)
upersonligt og interesseløst fortabende sig i det Objective, havde glemt
Tilegnelsen, glemt Opgaven, at existere i sin Viden; --- dels laa den i den
forvirrende Sammenblanding af de forskjellige Livsanskuelser, der var
begrundet i, at Verdsliggjørelsen af Christendommen og Forglemmelsen af
Personlighedens Lidenskab havde ført til Forglemmelsen af Christendommens
ideale Fordring. For at hæve Bedraget maatte da først disse Livsanskuelser
bringes frem i deres udprægede Eiendommelighed, og bringes frem, ikke som en
Gjenstand for Viden, thi saaledes vilde netop Tidens Retning bort fra
Existens og fra Existensens Afgjørelse begunstiges; men i en Form, der
hævder Existensens Betydning, der tvinger Læseren til \emph{ved sig selv
personligt} at erhverve Forstaaelsen. Han „gjøres opmærksom“, det Øvrige
overlades til ham selv; den \emph{indirecte} Meddelelse skjuler den
Meddelende for ham, gjør denne ligesom ukjendelig; han skal nu hjælpe sig
selv i Tilegnelsen, og i Tilegnelsen af det, der ved at fremstilles i
Existens kun \emph{indirecte} meddeles, skal han \emph{personligt erfare},
hvor den lavere Livsanskuelses Anstødssten ligger, hvor den støder an i
Lidenskab, hvor Existensen taber sit Grundlag, og hvor Fortvivlelsen fører
enten ud i Fortabelsen eller til det Høiere, til „at frelses uendeligt i
Religieusitetens Væsenlighed“ (jfr. En literair Anmeldelse. 1846. S. 108
ff.). Men idet de tidligere Livsstadier skulde fremstilles, ikke docerende,
som Gjenstand for Viden; men i existerende Individualiteter, maatte den
første Række af Skrifter \emph{blive pseudonym}. Pseudonymiteten havde „sin
væsenlige Grund i selve Frembringelsen, der for Replikens, for den
psychologisk varierede Individualitets Forskjelligheds Skyld digterisk
krævede den Hensynsløshed i Godt og Ondt, i Sønderknuselse og Overgivenhed,
i Fortvivlelse og Overmod, i Lidelse og Jubel osv., der kun er ideelt
begrændset af den psychologiske Conseqvens, hvilken ingen factisk virkelig
Person i Virkelighedens sædelige Begrændsning tør tillade sig, eller kan
ville tillade sig“ (Anhang til Afsl. Efterskrift). Pseudonymerne ere
saaledes \emph{digtede} Forfattere, ligesom dramatiske Personer, der
\emph{leve} i deres Anskuelser og udtale disse i hele deres ideale
Conseqvens.

I det første af de pseudonyme Skrifter: \emph{Enten---Eller,} et
Livsfragment, udgivet af Victor Eremita (1843) er nu de to Livsstadier, der
ligge forud for det Religieuse: det Æsthetiske og det Ethiske, fremstillet i
existerende Personligheder. „Det Æsthetiske er det i Mennesket, ved hvilket
det \emph{umiddelbart} er det, det er; den, der lever i og ved og af og for
dette, lever æsthetisk“; hans Løsen er \emph{Nydelsen}, hans Maal
Fortabelsen. Første Del fremstiller nu en saadan æsthetisk Individualitet,
omtumlet i Muligheder, som hans Phantasi fremmaner, hans Dialektik atter
tilintetgjør, splittet ad i de isolerede Øieblikkes Lidenskab, uden
Existensens Sammenhæng, uden Virkelighed, fangen af Fortvivlelsen. Den anden
Del fremstiller en enkelt Individualitet, der i Fortvivlelsens Opgiven af
det Æsthetiskes Umiddelbarhed har gjort Uendelighedens Bevægelse, hvorved
det Ethiske først kommer tilsyne, der i Fortvivlelsen har fundet sig selv
igjen i sin evige Gyldighed, \emph{valgt} sig selv, som i dette Valg er
bleven \emph{aabenbar} i det Almene, og ved dette har faaet Muligheden til
at give sit Liv Continuitet. Værkets Slutning, en opbyggelig Tale, fører hen
til Grændsen af det Religieuse, og ved at bringe det Opbyggelige ind,
accentueres Subjectivitetens, Tilegnelsens Betydning i Forholdet til
Sandheden.

Efter Enten---Eller fulgte nu tre mindre, pseudonyme Skrifter: \emph{Frygt
og Bæven, Gjentagelsen} (1843) og \emph{Begrebet Angest} (1844), der
successivt bringe de \emph{religieuse} Bestemmelser frem, medens
Enten---Eller gav de æsthetiske og ethiske. De fremstille i Existens
Collsionen mellem det Religieuse og det Ethiske, Prøvelsen, Kampen, i
% PRINTED AS IS: "Collsionen" for "Collisionen" — the i of the second
% syllable is simply absent; confirmed at 450% on the 400 dpi render.
hvilken Guds-Forholdet seirer, og endelig det Religieuses afgjørende
Bestemmelse: \emph{Synden}, der sætter Individet som uensartet med det
Ethiske, saaledes at det ikke i Fortvivlelsen \emph{ved sig selv} kan hævde
sig selv, vinde sig selv ud af Fortvivlelsen; men kun ved et Høiere, ved
Gudsforholdet, i Selvtilintetgjørelse. Samtidig med det sidste af disse
Skrifter udkom en lille Bog, fuld af sprudlende Vittighed og bidende Satire:
„\emph{Forord}, Morskabslæsning for enkelte Stænder efter Tid og
Leilighed“, der i sit sidste Afsnit, spøgende, men alvorligt, viser hen paa
den religieuse Ligelighed, paa Fordringen af, at det Høieste for Mennesket
maa være det for Alle Tilgængelige til Forskjel fra Talentets tilfældige
Differens.
% ---- COLUMN BREAK: col. 2 ends "…Talentets tilfældige Differens."
%      col. 3 opens "(„Denne Menneskeligheds og Menneske-Lighedens Tanke:…"
%      The folio "— 1179 —" stands below col. 2 at the foot of the leaf.
% (Column 3 opens FLUSH, not indented — measured at the native raster
% 2026-09-23 — so this is not a paragraph break.)
(„Denne Menneskeligheds og Menneske-Lighedens Tanke: christeligt er ethvert
Menneske (den Enkelte), ubetinget ethvert Menneske, Gud lige nær.“ (Forord
til To Taler ved Altergangen).
% PRINTED AS IS: the parenthesis opened before „Denne is never closed, and
% the inner quotation closes after "nær.“ while a second parenthesis opens
% at "(Forord". The brackets of this sentence do not balance on the page.

Det Religieuse var nu naaet i dets væsenlige Bestemmelse, i \emph{Syndens}
Bestemmelse. Denne var ikke udviklet docerende, „thi Synden har, ligesom
Troen, Paradoxet o. desl. ingen Plads i Systemet“, men bragt frem i dens
psychologiske Existens i Individet som \emph{Angest}. Nu var det sikkert, at
det Christelige ikke forvexledes med noget Lavere, og „\emph{Philosophiske
Smuler}“ (1844) fremsatte det nu i Form af en Hypothese, experimenterende.
\emph{Formen} betingede her Pseudonymitet, \emph{Gjenstanden} Navngivelse,
derfor findes K.s Navn paa dette Skrift som \emph{Udgiver}. Problemet i
dette Skrift er: „Kan der gives et historisk Udgangspunct for en evig
Bevidsthed; hvorledes kan et saadant interessere mere end historisk; kan man
bygge en evig Salighed paa en historisk Viden?“ Hypothesen i det er:
„\emph{hvis} der skal gaaes videre end det Sokratiske“, det vil sige end den
af Immanensens Tanke baarne Livsanskuelse, hvor Tiden, Endeligheden, bliver
det Forsvindende som evigt optaget i det Uendelige; hvor Evighedens
Afgjorthed gjør Existensen, hvor alvorligt den end tages, til en Spøg; hvor
Læreren kun bliver \emph{Anledningen} for den Lærende, fordi denne evigt
\emph{har} Sandheden, hvor selve \emph{Skyldens} Bevidsthed ikke naaer ud
over Døden, fordi den endelige Personlighed ikke naaer ud over denne.
„\emph{Skal} der altsaa gaaes videre“, saa maa andre afgjørende
Bestemmelser sættes, og nu kommer gjennem en Construction, der parodierer
Speculatione, de orthodox-christelige Bestemmelser frem i deres hele
% PRINTED AS IS: "Speculatione" for "Speculationen" — the final n is absent;
% confirmed at 250% on the 400 dpi render.
Simpelhed. Hypothesen bliver nu: „et nyt Organ; \emph{Troen,} en ny
Forudsætning: \emph{Syndsbevidstheden,} en ny Afgjørelse: \emph{Øieblikket,} en ny Lærer:
\emph{Guden i Tiden}“.
% UNSURE: in "Guden i Tiden" the one-letter word "i" cannot be measured —
% a single letter has no internal advance — so whether it belongs to the
% letterspaced run is undecidable at this scan's 200 ppi. Both flanking
% words are spaced; the run is marked whole.

Samtidig med disse pseudonyme Skrifter havde K. under sit eget Navn udgivet
6 smaa Samlinger af \emph{opbyggelige Taler} (1843---1844), der med
bestandigt Sigte paa Religionens Kategori: „\emph{den Enkelte}“\footnote{%
% The printed mark is *), set after the closing „ of „den Enkelte“.
„Det er christelig Heroisme --- at vove ganske at blive sig selv, et
\emph{enkelt} Menneske, dette bestemte \emph{enkelte} Menneske, ene ligeoverfor
Gud, ene i denne uhyre Anstrængelse og dette uhyre Ansvar“.
% The s of "Anstrængelse" has barely taken ink — only a faint spine survives,
% so the word reads as "Antrængelse" on the leaf. Absent ink, not another
% sort: transcribed as printed intent, per the standing rule.
Sygd. til Døden, Fortale. --- Det var til denne „\emph{Enkelte}, hvem jeg
med Glæde og Taknemmelighed kalder min Læser“ at Forordene til alle de
opbyggelige Taler henvendte sig. Om den Enkeltes Forhold til „Menighed“ har
K. udtalt sig udførligere i Indøvelse i Christendom S. 235 f. (ny Udgave),
jfr. Om min Forfattervirksomhed, S. 12. Anm.; Litter. Anmeld. S. 107 f.%
% This quotation DOES close: the mark after "min Læser" is a raised closing “
% (baseline -11.2 to -6.2 px, matching the closing “ after "Ansvar“", while an
% opening „ on the same leaf sits on and below the baseline). An earlier note
% claiming it unclosed was wrong; corrected at a third reading 2026-09-22.
} nærmede sig mere og mere hen imod det Christeligt-Religieuse, uden dog at
ville naae dette, idet de afsluttede i det Humoristiske.

Ved Slutningen af \emph{Philosophiske Smuler} var der antydet et andet
Skrift, i hvilket „Problemet skulde ifores historisk Costume“. Men før dette
% UNSURE: "ifores" — the sort shows no ø-slash at 400 dpi, and the word is
% not letterspaced, so it is not the page's bold-face habit. Read as printed;
% the alternative rejected is "iføres".
udkom endnu „\emph{Stadier paa Livets Vei}“ (1845), i hvilket de tre
Existensstadier: det Æsthetiske, Ethiske og Religieuse bestemtere
udprægedes, og det religieuse Livs Grundbestemmelse: \emph{Lidelsen}, skarpt
accenteredes. Aaret efter udkom da „\emph{Afsluttende uvidenskabelig
Efterskrift til philosophiske Smuler}, ligesom „Smulerne“ med K.s Navn som
% PRINTED AS IS: the letterspaced title is opened with „ but never closed
% before the comma; "ligesom „Smulerne“" then opens a second pair inside it.
\emph{Udgiver}. Her stilles nu afgjørende Problemet om Subjectets Forhold
til Christendommen, eller: det, \emph{at blive Christen}; Subjectiviteten
som Sandheden for den Existerende hævdes mod Objectivitetstheorierne,
Existensens Betydning hævdes mod Speculationens Seen bort fra Existens,
indirecte paapeges det Punct, hvor de tidligere Livsanskuelser maae strande;
Existenssphærerne sondres bestemtere ud fra hverandre; Ironien bestemmes som
Confiniet mellem det Æsthetiske og det Ethiske, Humor som Confiniet mellem
det Ethiske og Religieuse; den religieuse Sphære deles i to: den
Religieusitet, der hviler i Immanensen, og den christelige, hvis
Grundbestemmelse, „i hvilken alle de andre christelige Bestemmelser ligge,
og af hvilken de conseqvent lade sig udlede“, er den, „at i \emph{Tiden}
afgjøres Spørgsmaalet om Individets evige Salighed og af jøres ved Forholdet
% PRINTED AS IS: "af jøres" for "afgjøres" — the g is absent and the space
% it left stands open between "af" and "jøres"; confirmed at 450% on the
% 400 dpi render.
til Christendommen som noget Historisk“, i hvilken Grundbestemmelse
Forargelsens Mulighed ligger.

Med dette Skrift var der afgjørende naaet \emph{hen til} Maalet, til det
Christelige; dette var bragt frem i sin absolute Forskjel fra de tidligere
Stadier, ikke som en Gjenstand for Viden, men som Troens Gjenstand:
Paradoxet, Videns Fortvivlelse, Forargelsens Anstødssten. Det Reli%
% Col. 3, and printed p. 1179, break off here in mid-word: "Det Reli-".
% The printed hyphen is resolved and the word joined across the page turn,
% "Reli\opage{1180}gieuse", as elsewhere in this collection; the trailing %
% on the line above and the absence of any blank line below keep it one word.
% The word is completed at the head of p. 1180 ("gieuse"), so the next
% fragment's page marker falls INSIDE this word: its opening marker belongs
% between "Reli" and "gieuse", not before "Det". (The marker is deliberately
% not spelled out here, so that a grep for page markers in THIS fragment
% returns 1179 and nothing else.)
% --- p. 1180 ---
% Fædrelandet nr. 281, 1. December 1855, printed p. 1180 = PDF page 3.
% Folio "1180" verified at the head of the leaf (centred, between rules).
% The page is set in THREE columns of ANTIQUA (roman), not Fraktur: the
% body of this newspaper is roman throughout, so emphasis is by
% letterspacing (-> \emph{}). There is NO italic anywhere in the article:
% "tilbage i Enfold" was first read as sloped, and a third reading at the
% native 200 ppi raster settled it as ordinary roman (double-storey a,
% binocular g, stems within 4 degrees of vertical, advance tighter than its
% roman neighbours). Corrected 2026-09-22.
% The article continues mid-sentence from p. 1179 at the head of col. 1
% and ENDS in col. 3, signed "-r." (see the note at the end).
\opage{1180}gieuse havde været Maalet, det havde „helt og holdent været sat ind i Reflexion, men saaledes, at det helt og holdent toges ud af Reflexion tilbage i Enfold“, den tilbagelagte Vei havde været: „\emph{at naae, at komme til} Enfoldighed“, Opgaven „ikke at reflectere sig ind i Christendom; men at reflectere sig ud af Andet, og blive, enfoldigere og enfoldigere, Christen“.
% UNSURE (emphasis): the letterspacing of „at naae, at komme til stops at
% the line end; "Enfoldighed" at the head of the next line is set with
% ordinary spacing (checked at 900 dpi against the letterspaced
% "ninger," of Gjerninger on the same page). Emphasised only as far as
% the print letterspaces it.

Efter Rækken af de pseudonyme Skrifter fulgte nu Rækken af de \emph{opbyggelige}, først „\emph{Opbyggelige Taler i forskjellig Aand}“ (1847), i hvilke, som Titelen antyder, det Opbyggelige fra de tidligere Stadier var optaget, saa at kun den sidste Tredi del gav det eiendommeligt Christelige;
% PRINTED AS IS: "Tredi del", set as two words with a capital T
% (verified at 800 dpi); not normalised to "Trediedel".
derefter \emph{Kjærlighedens Gjerninger}, christelige Overveielser i Talers Form (1847) og \emph{Christelige Taler} (1848), til hvilke i de nærmest følgende Aar flere mindre Samlinger sluttede sig, bestandig b stemtere stræbende henimod
% PRINTED AS IS: "b stemtere" — the second sort of "bestemtere" has
% failed to print, leaving "b" + space + "stemtere" (verified at 800
% dpi). Transcribed as printed; intended reading "bestemtere".
hvad der er Christendommens Midtpunkt, for at, som K. dengang udtalte det, den hele Forfattervirksomhed maatte finde „sit afgjørende Hvilepunkt ved Alterets Fod, hvor Forfatteren, personligt sig sin Ufuldkommenhed og Skyld selv bedst bevidst, ingenlunde kalder sig et Sandhedsvidne, men kun en egen Art Digter og Tænker, der, „uden Myndighed“, intet Nyt har havt at bringe, men „villet læse de individuelle humane Existens-Forholds Urskrift, det Gamle, Bekjendte og fra Fædrene Overleverede igjennem endnu engang, om muligt paa en inderligere Maade“.“ (Forord til To Taler ved Altergang om Fredagen, 1851.)

De opbyggelige Skrifters Opgave var „med Guds Bistand at anvende Alt, for at faae det klart, hvad Christendommens Fordring i Sandhed er“, ikke at føre en Apologie for Christendommen, men „at fremstille den, hvad Sandhed er, som noget saa uendelig høit, at Apologien falder paa et andet Sted, bliver for os, at vi vove at kalde os Christne, eller forvandler sig til en bodfærdig Bekjendelse, at vi takke Gud, naar vi blot tør ansee os selv for Christne.“ Men ved Siden af den uendelige Fordrings Strænghed glemtes ikke Naadens Uendelighed, kun maatte Fordringen „høres og hævdes først, høres og hævdes i sin hele Uendelighed“, før „Naaden“ kunde anbringes. „Det, jeg har villet, er at bidrage til, ved Hjælp af Tilstaaelser at bringe om muligt noget mere Sandhed ind i disse (i Retning af at være ethisk og ethisk-religieus Charakter, at forsage verdslig Klogskab, at ville lide for Sandhed o. desl.) ufuldkomne Existenser, som vi føre dem, hvilket dog altid er Noget, og i ethvert Tilfælde den første Betingelse for at komme til at existere dygtigere. Det, jeg har villet forhindre er, at man da ikke, indskrænkende sig til og nøiende sig med det Lettere og Lavere, derpaa gaaer videre, afskaffer det Høiere, --- gaaer videre, sætter det Lavere i det Høieres Sted, --- gaaer videre, gjør det Høiere til Phantasteri og latterlig Overdrivelse. det Lavere til Visdom og sand Alvor, at man da ikke i „Christenhed“ existentielt tager Luther og Betydningen af Luthers Liv forfængeligt: dette har jeg villet bidrage til om muligt at forhindre“ (Om min Forfattervirksomhed S. 18).
% UNSURE: "afskaffer det Høiere" — the capital H is under-inked in this
% copy and prints almost as "løiere"; read as Høiere from the identical
% sort in "det Høieres Sted" on the next line (800 dpi).
% Overdrivelse. — the mark after "Overdrivelse" was first read as a comma.
% Two later readings at the native 200 ppi raster make it a FULL STOP: it is
% round, 6 px, ending +0.9 px below the baseline, matching "S. 18)." three
% lines below, while the comma of "Alvor," on the same line is 9 px with a
% tail running +4.5 px below. Corrected 2026-09-22. The next word is printed
% lowercase ("det") after the stop; transcribed as printed, not emended.

Christendommens Fordring i dens „Strænghed“ gjordes gjældende i to Skrifter, der udkom pseudonymt, med Forfatterens Navn som Udgiver. I det Første: \emph{Sygdommen til Døden}, en christelig psychologisk Udvikling til Opbyggelse og Opvækkelse (1849) er Sygdommen til Døden ɔ: \emph{Fortvivlelse} opfattet som Synden, denne som Modsætning til \emph{Tro}, hvis Gjenstand er Paradoxet, det for Tanken Absurde; \emph{Forargelsens} Mulighed fastholdt som det dialektiske Moment i alt det Christelige. I det andet Skrift, det, der betegner det afgjørende Vendepunkt i K.s Forhold til den bestaaende officielle Kirke: \emph{Indøvelse i Christendom} (udgivet 1850, skrevet 1848) blev „Fordringen til det at være Christen tvunget op til et Idealitetens Høieste“ i Retning af \emph{Samtidigheds-}Forhold til Forbilledet, af at gjøre Forbilledets Liv nærværende i vort, af Lidelse, Selvtilintetgjørelse, af Tro i Modstrid mod Forstanden under Forargelsens uadskillelige Mulighed.
% Emphasis: the letterspacing of "Samtidigheds-" stops at the line end;
% "Forhold" at the head of the next line is set normally.
Denne Bog var et Spørgsmaal, der blev stillet til det kirkeligt Bestaaende. I Opfattelsen af Lidelsen som Religieusitetens uadskillelige Særkjende, i Opfattelsen af Christendommen som Uensartethed med Verden, som Fordring til Efterfølgelse af Forbilledet, af „Omdannelse af Existensen i Lighed med Gud“ var der udtalt en Dom over en „Christenhed“, der har gjort sig det hyggeligt i Verden, fundet Hvile i Verdsligheden, over en Geistlighed, der, beskyttet af Staten, i Velværens \emph{Nydelse} forkynder \emph{Lidelsens} Religion. Hvad der sildigere blev udtalt i „Øieblikket“, var her ikke blot antydet, men i alt Væsenligt udtrykkeligt udtalt (jvfr. især S. 41, 75, 123, 136, 209, 228, 241, 243 f., 245, 268 i den nye Udgave).
% Emphasis: "Lidelsens" is letterspaced, "Religion" is not (800 dpi).
% "blev udtalt i „Øieblikket“" — the "i" is an x-height stem on the baseline
% with its dot BELOW the baseline, i.e. a TURNED SORT, not a battered letter
% and not an exclamation mark (which would stand ascender-high with its point
% on the baseline). Reading "i"; the turned sort is the printer's.
Men, utilbøielig som K. var til at virke henimod en Forandring i det Ydre, gjorde han endnu ved den Fortale, han føiede til Bogen, et Forsøg paa at give det Bestaaende et Forsvarsmiddel. „Min Tanke var: skal dette Bestaaende kunne forsvares, er dette den eneste Maade: ved digterisk (derfor af en Pseudonym) at anbringe Dommen over det, ved saa at trække i anden Potens paa „Naaden“, saa det blev Christendom, ikke blot ved „Naaden at finde Tilgivelse for det Forbigangne, men ved Naaden et Slags Aflad fra den egenlige Christi Efterfølgelse og den egenlige Anstrængelse ved det at være Christen.
% PRINTED AS IS: the inner quotation „Naaden at finde Tilgivelse ... is
% opened and never closed; quote balance drifts by +1 here.
Paa den Maade kommer der dog Sandhed ind i det Bestaaende; det forsvarer sig ved at dømme sig selv, det vedgaaer den christelige Fordring, gjør Tilstaaelse sig selv betræffende om sin Afstand, og uden at kunne kaldes en Stræ%
% COLUMN BREAK: col. 1 ends with the broken word "Stræ-" and col. 2
% opens with "ben i Retning af at komme Fordringen nærmere".
ben i Retning af at komme Fordringen nærmere, men henflyer til Naaden „ogsaa i Forhold til den Brug, man gjør af Naaden“. („Fædrelandet“ for 16de Mai 1855, jvfr. Udgiverens Forord til \emph{Indøvelse i Christendom}, og: Om min Forfattervirksomhed, S. 17, 19, Anm.)

% TRANSCRIBER'S ERROR CORRECTED 2026-09-23: "almægtig" -> "afmægtig". The
% second letter carries the hook of an f at the top of its stem, matching
% the f of "fordømme" on the next line; the l of "valgte" on the same line
% has none. Found because the translation (made from Koch) read "impotent".
Spørgsmaalet var nu stillet til den officielle Kirke, nærmest til dens Overhoved, afdøde Biskop Mynster; i Dommen over den vilde Kirkens Overhoved udtale sin og Kirkens egen Dom. „Hvis der var Kraft i ham, maatte han gjøre Et af To; \emph{enten} afgjort erklære sig for den Bog, vove at gaae med den, lade den være Forsvaret, der afværger, hvad Bogen digterisk indeholder, Sigtelsen mod hele den officielle Christendom, at den er en Øienforblændelse, „ikke en sur Sild værd“, \emph{eller} saa afgjort som muligt kaste sig imod den, stemple den som et gudsbespotteligt og vanhelligt Forsøg, og erklære den officielle Christendom for den sande Christendom“. Ingen af Delene skete: „afmægtig valgte han Intet at gjøre, uden høist at fordømme den i Dagligstuen; men ikke engang i privat Samtale med mig, uagtet jeg bad ham derom, efter at det med hans Samtykke var mig forebragt, hvad han dømte i Dagligstuen.“ („Fædrel.“ paa anf. Sted.)

Saaledes stod Sagen hen medens Mynster levede. Kun den Sky for at paaskynde en Afgjørelse i det Ydre, der var begrundet ligesaa meget i K.s Opfattelse af det Ethiskes og Religieuses Væsen, som i hans Opfattelse af den Udvikling, gjennem hvilken Tiden maatte helbredes (jfr. den mærkelige Slutning paa „En literair Anmeldelse.“ 1846. S. 107 ff. og Artiklen i „Fædrelandet“ for 31te Januar 1851), hans fra Barndommens Paavirkninger bevarede Pietet mod Mynster, og vel ogsaa Haabet om, mod al Sandsynlighed, at see den ham „saa uendelig kjære Tanke“ blive til Virkelighed, at Mynster endnu kunde sige: „dette er i Sandhed Christendom, saaledes forstaaer jeg selv Christendommen i stille Inderlighed, saa at Bogen blev en Forklarelse af Biskop Mynsters Christendoms Forkyndelse“, kunde endnu holde ham tilbage fra det afgjørende Skridt. Da døde Mynster og ved hans Grav forkyndte hans Eftermand ham som Sandhedsvidne, et af de rette Sandhedsvidner. Med denne Dom var Svaret givet paa det Spørgsmaal, der var stillet ved „Indøvelse i Christendom“. Den officielle Kirke havde nu erklæret, at dens Christendomsforkyndelse, „der fortonede og fortaug det afgjørende Christelige“, var den sande, at Christendommen i dens ideale Fordring, det nye Testamentes Christendom, var Overdrivelse, Overspændthed. Nu maatte Indsigelsen gjøres mod denne Kirke, nu først drev Modsætningens Nødvendighed til at føre Sagen til det Yderste (jfr. „Dette skal siges.“ S. 10). Dog skete det ikke overilet; først næsten et Aar efter Mynsters Død og hans Eftermands Kanoniseren af ham, blev Indsigelsen offenliggjort (i „Fædrel.“ for 18de Decbr. 1854). Hvad der siden den Tid er skeet, er endnu i frisk Minde. I Artiklerne i „Fædrelandet“ og i „Øieblikket“ kæmpedes der for Christendom imod Christenhed og dennes \emph{tause} Repræsentant, Geistligheden. Det blev nu udtalt: „Hvo Du end er, hvilket Dit Liv forresten er, min Ven; --- ved, (hvis Du ellers deltager i den) at lade være at deltage i den offenlige Gudsdyrkelse som den nu er (med Paastand paa at være det nye Testamentes Christendom) har Du bestandig een og en stor Skyld mindre, Du deltager ikke i at holde Gud for Nar ved at kalde det nye Testamentes Christendom hvad der ikke er det nye Testamentes Christendom.“ („Dette skal siges, saa være det da sagt.“ S. 3). Der stilledes den Fordring til Staten, „at tage det Sandsebedrag bort med en Kirke, der i Verdslighed var smeltet sammen med det Verdslige, med en Geistlighed, „hvis Tilværelse, christeligt, er en Usandhed“, „hvis christelige Alvor finder sit Udtryk i, at „Sandhedsvidner“ gjøre, --- hvilken satirisk Selvmodsigelse! --- Carriere og Lykke i Verden ved om Søndagen at skildre hvorledes Sandheden maa lide i denne Verden!“, hvis hele Stilling tvinger den til at skjule for Menigheden, hvad Christendom er, og saaledes, saavidt det er afhængigt af den, umuliggjøre denne. Alle de Modsigelser, der ligge i Geistlighedens Stilling, alt det Blændværk, den har bragt ind i „Christenheden“, blottedes med bidende Satire: „Den Almægtige veed bedst, hvorledes der skal slaaes, saa det føles, at Latteren, betjent i Frygt og Bæven, maa være Svøben --- derfor bruges jeg.“ Og mod den usle Halvhed, der af Christendommen blot har Navnet tilbage, mod den Indbildning, „at det at være Christen, er Noget, som Enhver uden videre er, Noget, man slipper lettere til end den allerubetydeligste Færdighed“, stilledes Idealet alvorligt og dømmende. Det er nu 2 Maaneder siden, at det sidste Numer af „Øieblikket“ udkom; faa Dage efter fængslede Sygdom Forfatteren til det Sygeleie, hvorfra han ikke mere reiste sig.
% PRINTED AS IS: in this paragraph the quotation opened at
% „at tage det Sandsebedrag bort med en Kirke ... is never closed:
% the inner quotations close in turn, and the text runs on into
% „... umuliggjøre denne. Alle de Modsigelser ...“ without a
% closing mark (verified at 800 dpi).
% PRINTED AS IS: "Numer", one m, as the compositor sets it (800 dpi).
Hans Virksomhed er nu afsluttet, og naar vi see tilbage over den, ville vi sande hans Ord, at han „som Forfatter kun har villet Eet. At gjøre opmærksom paa det Christelige, var hans Livs Opgave, det var Opgaven i den første Række af hans Skrifter, der arbeidede mod det Christeliges Sammenblanden med hvad der ikke er christeligt, mod Forvexlingen af en Viden om Christendommen og en Forvandlen af Existensen ved Christendommen; det var Opgaven i hans opbyggelige Skrifter, der, med dyb Kjendskab til det menneskelige Hjertes skjulte Lidenskaber, forstode at opspore ethvert Selvbedrag, at spærre enhver Udvei, ad hvilken det kunde udflyes at \emph{see} det Christelige; det var Opgaven i hans Kamp mod den Kirke, der fortaug det Christelige, „gjorde Christendommen til Poesi,
% COLUMN BREAK: col. 2 ends with "til Poesi," and col. 3 opens with
% "til Noget, hvortil man kun forholder sig gjennem Indbildningskraften".
til Noget, hvortil man kun forholder sig gjennem Indbildningskraften“, afskaffede Christi Efterfølgelse i Lidelse, i Afdøen fra Verden, der „kastede Forbilledet bort, og lod det, man selv er, Middelmaadighed, gjælde for Idealet.“ Og hvad han i sit tidligt endte Liv har kæmpet for, med en Kraft i Villien, der ikke lod sig bøie af Legemets Svaghed, og med en Tro, der ikke forsagede, det har ikke været uden Frugt. Hans tidligere Skrifter have i Stilhed virket i en stedse større Kreds; hans „Øieblikke“ have hos Tusinder aabnet Øiet og vakt alvorlige Tanker, der ikke ville glemmes nu da han er død; men finde nye Næring i de Skrifter, han har efterladt, i det Liv, han har vakt i saa Manges Sjæle. \hfill ---r.
% PRINTED AS IS: the long quotation „som Forfatter kun har villet Eet.
% is opened and closed only once, at "gjælde for Idealet.“", although
% two further quotations are opened inside it; the quote balance on this
% page therefore drifts by +1 here as well (+2 for the page).
%
% END OF ARTICLE. The last words are "... i det Liv, han har vakt i saa
% Manges Sjæle." in COLUMN 3, about a sixth of the way down the column,
% followed on the SAME line, set flush to the right-hand edge of the
% column, by the signature mark "---r." (an em-dash, r, full stop).
% UNSURE (signature): the mark is over-inked in this copy — at 1200 dpi
% the "r" is a solid blob — but its width, the dash before it and the
% point after it agree with "-r", the mark the article is signed with.
%
% QUOTE BALANCE: the body of this fragment has 51 „ against 47 “.
% The drift of +4 is the print's own: three quotations are opened and
% never closed („Naaden at finde Tilgivelse ..., „at tage det
% Sandsebedrag ..., and „som Forfatter kun har villet Eet ..., which is
% closed only once, after "gjælde for Idealet.“", although two further
% quotations are opened inside it), and the missing closer of the first
% of them leaves „Min Tanke var: ... without one as well. Nothing has
% been supplied.
%
% WHAT FOLLOWS ON THE PAGE: a short centred rule closes the article;
% then unrelated newspaper matter in col. 3 — a political note ("--- Grev
% Holstein-Holsteinborg, Baron Blixen-Finecke og Baron Constant
% Dirckinck-Holmfeldt ..." on a meeting of Rigsraadsvælgere at Roskilde),
% an obituary note of Dr. med. R. Hübertz, a weekly note of deaths and
% illnesses reported by the Copenhagen physicians, and then a long
% advertisement, set in a smaller antiqua, for a Dampkjøkken.
% WHAT PRECEDES IT: nothing — the article occupies the head of col. 1,
% continuing mid-sentence ("...gieuse havde været Maalet") from p. 1179.

\end{document}
